
\chapter*{RINGKASAN}
\addcontentsline{toc}{section}{
    \bfseries\protect\numberline{}RINGKASAN
}

Laporan kemajuan ini mendokumentasikan pelaksanaan program magang industri selama satu semester pertama (Januari – April 2026) dari total durasi satu tahun di PT Inti Utama Solusindo, unit bisnis pendukung teknologi informasi bagi Pharos Group. Sebagai Front-End Developer, fokus utama pekerjaan melibatkan pemeliharaan dan migrasi sistem pada aplikasi ekosistem kesehatan dan farmasi. Aktivitas teknis mencakup pembaruan versi Flutter pada aplikasi HealthyOne dari versi satu ke versi dua, perbaikan bug, serta inisiasi migrasi arsitektural dari Flutter ke React untuk meningkatkan skalabilitas. Dalam proses pengembangan, diterapkan prinsip Clean Code dan SOLID guna memastikan kode yang dihasilkan bersifat maintainable. Manajemen proyek dilakukan secara profesional menggunakan OpenProject dengan sistem pemantauan kinerja berbasis poin (KPI). Hasil awal menunjukkan peningkatan pemahaman teknis mengenai integrasi antarmuka dengan layanan Back-End serta pentingnya peran Quality Assurance dalam siklus hidup pengembangan perangkat lunak.

Kata Kunci: MBKM, Front-End Developer, Flutter, React, SOLID.